\chapter{Teori Grup}\label{ch:b3:groups}

Pada jilid Tahun ke-2 grup diperkenalkan sebagai alat pembukuan: teorema
Lagrange, grup siklik, grup simetri beserta tandanya. Bab ini
mengubah teori grup menjadi sebuah \emph{metode}. Mesinnya adalah
gagasan grup yang \emph{beraksi} pada sebuah himpunan: menghitung
\omterm{def:b3:groups:action}{orbit} dan titik tetap melahirkan \omterm{cor:b3:groups:classeq}{persamaan kelas}, teorema Cauchy dan
ketiga teorema Sylow --- asas lokal-ke-global yang menjadi tulang
punggung teori grup berhingga. Setelah itu kita belajar merakit grup
(\omterm{prop:b3:groups:direct}{hasil kali langsung} dan semilangsung) dan membongkarnya (\omterm{thm:b3:groups:jordanholder}{deret
komposisi}, \omterm{def:b3:groups:derived}{grup terselesaikan}), lalu membuktikan teorema yang, pada
\cref{ch:b3:galois}, akan menutup pertanyaan berumur tiga abad
tentang persamaan polinomial: \omterm{thm:b3:groups:ansimple}{grup alternasi} $A_n$ \omterm{def:b3:groups:simple}{sederhana} untuk $n
\geq 5$.

\section{Grup kuosien dan teorema isomorfisma}

Sepanjang bab ini $G$ adalah grup yang ditulis secara multiplikatif,
dengan unsur identitas $e$. Ingat kembali dari jilid Tahun ke-2:
subgrup, koset $gH$, teorema Lagrange ($\abs G = [G:H]\,\abs H$ untuk
$G$ berhingga), orde sebuah unsur, grup siklik, dan grup simetri $S_n$
beserta morfisma tandanya $\varepsilon \colon S_n \to \{\pm1\}$.

\begin{definition}\label{def:b3:groups:normal}
Sebuah subgrup $N$ dari $G$ disebut \emph{normal}\index{subgrup normal}
(ditulis $N \trianglelefteq G$) apabila $gNg^{-1} = N$ untuk setiap $g
\in G$ --- ekuivalen dengan: koset kiri dan koset kanan berimpit, $gN =
Ng$ untuk setiap $g$.
\end{definition}

\begin{theorem}[Grup kuosien]\label{thm:b3:groups:quotient}
Misalkan $N \trianglelefteq G$. Himpunan koset $G/N$, dengan
perkalian $(gN)(hN) = ghN$, adalah grup yang terdefinisi dengan baik,
yaitu \emph{grup kuosien}\index{grup kuosien}, dan
\emph{proyeksi kanonik} $\pi \colon G \to G/N$, $g \mapsto gN$,
merupakan morfisma surjektif dengan kernel $N$. Sebaliknya, setiap
kernel morfisma grup bersifat \omterm{def:b3:groups:normal}{normal}: \omterm{def:b3:groups:normal}{subgrup normal} tidak lain
adalah kernel.
\end{theorem}

\begin{proof}
Keterdefinisian dengan baiklah yang menjadi inti persoalan. Jika $gN = g'N$ dan $hN = h'N$,
tulis $g' = gn$, $h' = hm$ dengan $n, m \in N$. Maka $g'h' = gnhm =
gh\,(h^{-1}nh)\,m \in ghN$ karena $h^{-1}nh \in N$ menurut kenormalan:
hasil kali koset tidak bergantung pada wakil yang dipilih.
Keasosiatifan, identitas $eN = N$ dan invers $(gN)^{-1} = g^{-1}N$
diwarisi dari $G$. Jelas bahwa $\pi$ morfisma surjektif dan
$\pi(g) = N \iff g \in N$.

Jika $f \colon G \to H$ sebuah morfisma dan $k \in \ker f$, maka
$f(gkg^{-1}) = f(g)f(k)f(g)^{-1} = e$: kernel bersifat \omterm{def:b3:groups:normal}{normal}.
\end{proof}

\begin{theorem}[Sifat universal; teorema isomorfisma pertama]
\label{thm:b3:groups:firstiso}
Misalkan $f \colon G \to H$ sebuah morfisma dan $N \trianglelefteq G$
dengan $N \subseteq \ker f$. Terdapat tepat satu morfisma $\bar f \colon G/N
\to H$ yang memenuhi $f = \bar f \circ \pi$. Khususnya, dengan mengambil $N = \ker
f$:\index{teorema isomorfisma}
\[
G/\ker f \;\xrightarrow{\;\sim\;}\; \operatorname{im} f,
\qquad gN \mapsto f(g).
\]
\end{theorem}

\begin{proof}
Ketunggalan: $\bar f(gN)$ haruslah $f(g)$. Keberadaan: jika $gN = g'N$
maka $g^{-1}g' \in N \subseteq \ker f$, sehingga $f(g) = f(g')$ dan $\bar
f(gN) = f(g)$ terdefinisi dengan baik; ini morfisma karena $f$ morfisma. Untuk
$N = \ker f$: $\bar f$ injektif, sebab $\bar f(gN) = e$ berarti
$g \in \ker f$, yaitu $gN = N$; sedangkan petanya sama dengan peta $f$.
\end{proof}

\begin{theorem}[Teorema isomorfisma kedua dan ketiga]
\label{thm:b3:groups:secondiso}
Misalkan $H \leq G$ dan $N \trianglelefteq G$.
\begin{enumerate}
  \item $HN = \{hn : h \in H,\, n \in N\}$ adalah subgrup, $N
        \trianglelefteq HN$, $H \cap N \trianglelefteq H$, dan
        \[
        H/(H \cap N) \;\cong\; HN/N .
        \]
  \item Jika di samping itu $N \subseteq K \trianglelefteq G$, maka $K/N
        \trianglelefteq G/N$ dan $(G/N)\big/(K/N) \cong G/K$.
\end{enumerate}
\end{theorem}

\begin{proof}
(1) $HN$ adalah subgrup: $(hn)(h'n') = hh'\,(h'^{-1}nh')n' \in HN$
dan $(hn)^{-1} = h^{-1}(hn^{-1}h^{-1}) \in HN$, dengan memakai kenormalan
$N$. Susunlah $H \hookrightarrow HN \xrightarrow{\pi} HN/N$: morfisma
ini surjektif ($hnN = hN$) dengan kernel $\{h \in H : h \in
N\} = H \cap N$; lalu terapkan \cref{thm:b3:groups:firstiso}.

(2) Proyeksi $G/N \to G/K$, $gN \mapsto gK$, terdefinisi dengan baik
($N \subseteq K$), surjektif, dan berkernel $K/N$; terapkan
\cref{thm:b3:groups:firstiso} sekali lagi.
\end{proof}

\begin{theorem}[Teorema korespondensi]\label{thm:b3:groups:correspondence}
Misalkan $N \trianglelefteq G$. Pemetaan $H \mapsto H/N$ adalah bijeksi
antara subgrup $G$ yang memuat $N$ dan subgrup $G/N$; bijeksi ini
mempertahankan inklusi, indeks dan kenormalan (dalam kedua arah).
\end{theorem}

\begin{proof}
Inversnya adalah $\bar H \mapsto \pi^{-1}(\bar H)$. Kedua pemetaan mengirim
subgrup ke subgrup, dan keduanya saling invers: $\pi^{-1}(H/N) =
HN = H$ karena $N \subseteq H$, dan $\pi(\pi^{-1}(\bar H)) = \bar H$
karena $\pi$ surjektif. Inklusi jelas terpelihara;
$[G:H] = [G/N : H/N]$ sebab $gH \mapsto (gN)(H/N)$ adalah bijeksi
yang terdefinisi dengan baik antara ruang koset; dan $gHg^{-1} = H$ berlaku untuk
setiap $g$ jika dan hanya jika $(gN)(H/N)(gN)^{-1} = H/N$ untuk setiap $gN$, sekali lagi
karena $\pi$ surjektif.
\end{proof}

\begin{example}\label{ex:b3:groups:firstiso}
$\varepsilon \colon S_n \to \{\pm 1\}$ memberikan $S_n/A_n \cong
\{\pm1\}$; $\det \colon GL_n(K) \to K^\times$ memberikan
$GL_n(K)/SL_n(K) \cong K^\times$; $t \mapsto \eu^{2\iu\pi t}$ memberikan
$\R/\Z \cong \mathbb U$, yaitu grup lingkaran. Teorema isomorfisma
pertama adalah cara kuosien \emph{dihitung} dalam praktik: carilah
sebuah surjeksi yang kernelnya tepat seperti yang diinginkan.
\end{example}

\begin{method}\label{met:b3:groups:normal}
Untuk membuktikan $N \trianglelefteq G$, berurutan dari yang paling elegan:
tampilkan $N$ sebagai kernel sebuah morfisma yang terdefinisi pada $G$; periksa
$gNg^{-1} \subseteq N$ untuk setiap $g$ (ini sudah cukup: menerapkannya pada
$g^{-1}$ lalu mengonjugasikan memberikan inklusi sebaliknya); pastikan bahwa
$N$ merupakan gabungan \omterm{ex:b3:groups:actions}{kelas konjugasi}; atau amati bahwa $[G:N] = 2$
(maka $gN = Ng$ terpaksa berlaku --- \cref{exo:b3:groups:1}).
\end{method}

\section{Aksi grup}

\begin{definition}\label{def:b3:groups:action}
Sebuah \emph{aksi}\index{aksi grup} dari $G$ pada himpunan $X$ adalah
morfisma $\varphi \colon G \to \mathfrak{S}(X)$ ke dalam grup semua
bijeksi $X$; ditulis $g \cdot x$ untuk $\varphi(g)(x)$.
Secara ekuivalen: pemetaan $G \times X \to X$ dengan $e \cdot x = x$ dan $g
\cdot (h \cdot x) = (gh) \cdot x$. Adapun \emph{orbit}\index{orbit} dari
$x$ adalah $\mathcal O_x = \{g \cdot x : g \in G\}$,
\emph{stabilisator}\index{stabilisator} miliknya adalah subgrup $G_x = \{g :
g\cdot x = x\}$, dan $X^G = \{x : \forall g,\ g \cdot x = x\}$ adalah
himpunan \emph{titik tetap}. Aksi disebut \emph{transitif} bila
orbitnya tepat satu, \emph{setia} bila $\varphi$
injektif, \emph{bebas} bila semua stabilisatornya trivial.
\end{definition}

\begin{example}\label{ex:b3:groups:actions}
Lima \omterm{def:b3:groups:action}{aksi} ini menggerakkan seluruh teori grup berhingga:
\begin{enumerate}
  \item $G$ pada dirinya sendiri melalui \emph{translasi kiri} $g \cdot x = gx$:
        bebas dan transitif.
  \item $G$ pada dirinya sendiri melalui \emph{konjugasi} $g \cdot x = gxg^{-1}$:
        \omterm{def:b3:groups:action}{orbitnya} adalah \emph{kelas konjugasi}\index{kelas
        konjugasi}, \omterm{def:b3:groups:action}{stabilisatornya} \emph{sentralisator}
        \index{sentralisator} $Z_G(x) = \{g : gx = xg\}$, titik
        tetapnya \emph{pusat}\index{pusat} $Z(G)$.
  \item $G$ pada ruang koset $G/H$ melalui $g \cdot xH = gxH$:
        transitif, dan \omterm{def:b3:groups:action}{stabilisator} koset $H$ adalah $H$ sendiri.
        Setiap \omterm{def:b3:groups:action}{aksi} transitif berbentuk demikian
        (\cref{exo:b3:groups:8}).
  \item $G$ pada himpunan subgrupnya melalui konjugasi: \omterm{def:b3:groups:action}{stabilisator}
        $H$ adalah \emph{normalisator}\index{normalisator} $N_G(H) =
        \{g : gHg^{-1} = H\}$, yakni subgrup terbesar dari $G$ yang di dalamnya
        $H$ \omterm{def:b3:groups:normal}{normal}.
  \item $S_n$ pada $\intint{1}{n}$: induk dari segala contoh.
\end{enumerate}
\end{example}

\begin{theorem}[Orbit--stabilisator]\label{thm:b3:groups:orbitstab}
Pemetaan $gG_x \mapsto g \cdot x$ adalah bijeksi yang terdefinisi dengan baik
$G/G_x \to \mathcal O_x$. Khususnya, untuk $G$
berhingga,\index{teorema orbit--stabilisator}
\[
\abs{\mathcal O_x} = [G : G_x] \quad\text{membagi } \abs G,
\]
dan, karena \omterm{def:b3:groups:action}{orbit} mempartisi $X$ (\omterm{def:b3:groups:action}{orbit} adalah kelas-kelas dari
relasi ekuivalensi $x \sim y \iff y \in \mathcal O_x$),
\[
\abs X = \sum_{i} \,[G : G_{x_i}]
\qquad (x_i\colon \text{satu titik per orbit}).
\]
\end{theorem}

\begin{proof}
Terdefinisi dengan baik dan injektif: $gG_x = hG_x \iff h^{-1}g \in G_x \iff
h^{-1}g \cdot x = x \iff g \cdot x = h \cdot x$; bacalah rantai itu
dalam dua arah. Kesurjektifan adalah definisi \omterm{def:b3:groups:action}{orbit} itu sendiri. Pernyataan
pencacahannya menyusul dari teorema Lagrange dan dari pemartisian
$X$ atas \omterm{def:b3:groups:action}{orbit}.
\end{proof}

\begin{corollary}[Persamaan kelas]\label{cor:b3:groups:classeq}
Untuk grup berhingga $G$, dengan memilih satu wakil $x_i$ pada setiap
\omterm{ex:b3:groups:actions}{kelas konjugasi} yang beranggota lebih dari satu:\index{persamaan kelas}
\[
\abs G = \abs{Z(G)} + \sum_i \,[G : Z_G(x_i)],
\qquad\text{setiap } [G : Z_G(x_i)] > 1 \text{ membagi } \abs G.
\]
\end{corollary}

\begin{proof}
Terapkan \cref{thm:b3:groups:orbitstab} pada \omterm{def:b3:groups:action}{aksi} konjugasi:
\omterm{def:b3:groups:action}{orbit} beranggota tunggal tepat merupakan unsur-unsur $Z(G)$.
\end{proof}

\begin{theorem}[Titik tetap grup-$p$]\label{thm:b3:groups:pfixed}
Misalkan $p$ bilangan prima. Sebuah \emph{grup-$p$}\index{grup-p@grup-$p$} adalah
grup berhingga yang ordenya berupa pangkat $p$. Jika suatu grup-$p$ $G$ beraksi
pada himpunan berhingga $X$, maka
\[
\abs{X^G} \equiv \abs X \pmod p .
\]
Akibatnya: grup-$p$ yang tak trivial mempunyai pusat yang tak trivial, dan
setiap grup berorde $p^2$ bersifat abelian.
\end{theorem}

\begin{proof}
Setiap \omterm{def:b3:groups:action}{orbit} berkardinalitas $[G:G_x]$, sebuah pangkat $p$; pangkat itu
bernilai $1$ tepat pada titik tetap dan selain itu habis dibagi $p$. Menjumlahkan
atas semua \omterm{def:b3:groups:action}{orbit} memberikan kekongruenan tersebut. Untuk pusatnya: \omterm{def:b3:groups:action}{aksi}
konjugasi $G$ pada dirinya sendiri mempunyai $X^G = Z(G)$, sehingga $\abs{Z(G)} \equiv
\abs G \equiv 0 \pmod p$, dan $Z(G) \ni e$ memaksa $\abs{Z(G)} \geq
p$. Untuk orde $p^2$: jika $Z(G) \neq G$ maka $\abs{Z(G)} = p$ dan
$G/Z(G)$ siklik berorde $p$, yang memaksa $G$ abelian
(\cref{exo:b3:groups:2}) --- kontradiksi.
\end{proof}

\begin{theorem}[Cauchy]\label{thm:b3:groups:cauchy}
Jika bilangan prima $p$ membagi $\abs G$, maka $G$ memuat sebuah unsur
berorde $p$.\index{teorema Cauchy}
\end{theorem}

\begin{proof}[Bukti (McKay)]
Misalkan $X = \{(g_1, \dots, g_p) \in G^p : g_1 g_2 \cdots g_p = e\}$.
Memilih $g_1, \dots, g_{p-1}$ secara bebas sudah menentukan $g_p$: $\abs X =
\abs G^{p-1}$, habis dibagi $p$. Grup siklik $\Z/p\Z$ beraksi pada
$X$ melalui pergeseran siklik $(g_1, \dots, g_p) \mapsto (g_2, \dots, g_p,
g_1)$ --- \omterm{def:b3:groups:action}{aksi} ini mempertahankan $X$, sebab $g_2 \cdots g_p g_1 =
g_1^{-1}(g_1 \cdots g_p)g_1 = e$. Menurut \cref{thm:b3:groups:pfixed},
$\abs{X^{\Z/p\Z}} \equiv \abs X \equiv 0 \pmod p$. Titik tetapnya adalah
tupel konstan $(g, \dots, g)$ dengan $g^p = e$; tupel $(e,
\dots, e)$ salah satunya, jadi paling sedikit ada $p$ buah, sehingga
paling sedikit ada satu $g \neq e$ dengan $g^p = e$: ordenya tepat $p$.
\end{proof}

\begin{theorem}[Cayley]\label{thm:b3:groups:cayley}
Setiap grup berorde $n$ terbenam di dalam $S_n$.\index{teorema Cayley}
\end{theorem}

\begin{proof}
Translasi kiri $\varphi \colon G \to \mathfrak S(G) \cong S_n$ adalah
morfisma; $\varphi(g) = \mathrm{id}$ memaksa $g = ge = e$: jadi ia
setia.
\end{proof}

\begin{method}\label{met:b3:groups:counting}
Pencacahan titik tetap adalah langkah pembuka universal dalam teori
grup berhingga. Untuk membuktikan bahwa sesuatu \emph{ada} (unsur pusat,
unsur berorde $p$, \omterm{def:b3:groups:normal}{subgrup normal}, titik tetap), pilihlah grup yang
tepat dan himpunan berhingga yang tepat, buatlah grup itu beraksi, lalu bandingkan
$\abs{X^G}$ dengan $\abs X$ modulo $p$, atau manfaatkan bahwa ukuran \omterm{def:b3:groups:action}{orbit} membagi
orde grup. Bukti
\cref{thm:b3:groups:pfixed,thm:b3:groups:cauchy} serta ketiga
teorema Sylow di bawah adalah lima variasi atas satu gagasan tunggal ini.
\end{method}

\section{Teorema Sylow}

Teorema Lagrange mengatakan bahwa orde sebuah subgrup membagi $\abs G$;
konversnya tidak berlaku ($A_4$, berorde $12$, tidak mempunyai subgrup berorde
$6$ --- \cref{exo:b3:groups:1}). Teorema Sylow menyelamatkan
konvers itu untuk pangkat bilangan prima, dan klausa pencacahannya adalah
alat umum paling tajam yang kita punya untuk menghasilkan \omterm{def:b3:groups:normal}{subgrup normal}.

\begin{definition}\label{def:b3:groups:sylow}
Tulis $\abs G = p^a m$ dengan $p \nmid m$. Sebuah \emph{subgrup
Sylow-$p$}\index{subgrup Sylow} dari $G$ adalah subgrup berorde
$p^a$ --- yaitu subgrup-$p$ berorde sebesar mungkin. Banyaknya
subgrup Sylow-$p$ dari $G$ dilambangkan dengan $n_p$.
\end{definition}

\begin{lemma}\label{lem:b3:groups:binom}
Jika $\abs G = p^a m$ dengan $p \nmid m$, maka $\dbinom{p^a m}{p^a}
\equiv m \pmod p$.
\end{lemma}

\begin{proof}
Di dalam $\mathbb F_p[X]$, mimpi mahasiswa baru $(1+X)^p = 1 + X^p$ (koefisien
$\binom pk$, $0<k<p$, habis dibagi $p$: bilangan $p$
membagi pembilang $\frac{p!}{k!(p-k)!}$ tetapi tidak membagi
penyebutnya) beriterasi menjadi $(1+X)^{p^a} = 1 + X^{p^a}$, sehingga
\[
(1+X)^{p^a m} = \bigl(1 + X^{p^a}\bigr)^m
= \sum_{k=0}^{m} \binom mk X^{k p^a} \quad\text{di } \mathbb
F_p[X].
\]
Samakan koefisien $X^{p^a}$: di ruas kiri $\binom{p^a
m}{p^a} \bmod p$, di ruas kanan $\binom m1 = m$.
\end{proof}

\begin{theorem}[Sylow I: keberadaan]\label{thm:b3:groups:sylow1}
Untuk setiap bilangan prima $p$, subgrup Sylow-$p$ dari $G$ selalu ada.
\end{theorem}

\begin{proof}[Bukti (Wielandt)]
Misalkan $\Omega$ himpunan semua \emph{himpunan bagian} $G$ yang berkardinalitas
$p^a$; $G$ beraksi pada $\Omega$ melalui translasi kiri $g \cdot S = gS$.
Menurut \cref{lem:b3:groups:binom}, $\abs\Omega = \binom{p^a m}{p^a}
\equiv m \not\equiv 0 \pmod p$, sehingga ada \omterm{def:b3:groups:action}{orbit} $\mathcal O_S$ yang
ukurannya relatif prima terhadap $p$ (seandainya $p$ membagi setiap ukuran \omterm{def:b3:groups:action}{orbit}, ia akan membagi
$\abs\Omega$). Misalkan $H = G_S$ \omterm{def:b3:groups:action}{stabilisator} dari $S$ semacam itu.
Karena $[G : H] = \abs{\mathcal O_S}$ relatif prima terhadap $p$ dan $p^a \mid
\abs G = [G:H]\,\abs H$, diperoleh $p^a \mid \abs H$. Sebaliknya, tetapkan
$s \in S$: pemetaan $H \to S$, $h \mapsto hs$, injektif dan
nilainya jatuh di $S$ karena $hS = S$; jadi $\abs H \leq \abs S = p^a$.
Dengan demikian $\abs H = p^a$.
\end{proof}

\begin{theorem}[Sylow II: dominasi dan konjugasi]\label{thm:b3:groups:sylow2}
Misalkan $P$ sebuah subgrup Sylow-$p$ dan $Q$ sembarang subgrup-$p$ dari $G$.
Maka $Q \subseteq gPg^{-1}$ untuk suatu $g \in G$. Khususnya semua
subgrup Sylow-$p$ saling konjugat, dan $P \trianglelefteq G \iff
n_p = 1$.
\end{theorem}

\begin{proof}
Biarkan $Q$ beraksi pada ruang koset $X = G/P$, yang berkardinalitas $m
\not\equiv 0 \pmod p$. Menurut \cref{thm:b3:groups:pfixed} yang diterapkan pada
grup-$p$ $Q$, $\abs{X^Q} \equiv m \not\equiv 0 \pmod p$: jadi ada
koset tetap $gP$, yakni $QgP = gP$, yakni $g^{-1}Qg \subseteq
P$. Jika $Q$ sendiri \omterm{def:b3:groups:sylow}{subgrup Sylow}, kesamaan orde mengubah
$Q \subseteq gPg^{-1}$ menjadi kesamaan. Akhirnya $P
\trianglelefteq G$ jika dan hanya jika konjugatnya $\{gPg^{-1}\}$ --- yang menurut
uraian di atas adalah \emph{semua} subgrup Sylow-$p$ --- menyusut menjadi
$\{P\}$.
\end{proof}

\begin{theorem}[Sylow III: pencacahan]\label{thm:b3:groups:sylow3}
$n_p \equiv 1 \pmod p$, dan $n_p = [G : N_G(P)]$, yang membagi
$m$.
\end{theorem}

\begin{proof}
Misalkan $\mathrm{Syl}_p$ himpunan semua subgrup Sylow-$p$; $G$ beraksi
padanya secara transitif melalui konjugasi (\cref{thm:b3:groups:sylow2}), dengan
\omterm{def:b3:groups:action}{stabilisator} $P$ berupa \omterm{ex:b3:groups:actions}{normalisator} $N_G(P) \supseteq P$: jadi $n_p = [G :
N_G(P)]$, dan $m = [G:P] = [G:N_G(P)]\,[N_G(P):P]$ menunjukkan $n_p \mid
m$.

Sekarang batasi \omterm{def:b3:groups:action}{aksi} itu pada $P$ dan hitunglah titik tetapnya. Jika $Q \in
\mathrm{Syl}_p$ tetap oleh $P$, maka $P \subseteq N_G(Q)$; baik
$P$ maupun $Q$ subgrup Sylow-$p$ dari grup $N_G(Q)$, sehingga keduanya
saling konjugat di dalamnya (\cref{thm:b3:groups:sylow2} diterapkan pada $N_G(Q)$);
tetapi $Q \trianglelefteq N_G(Q)$, jadi $Q$ satu-satunya konjugat dirinya di sana:
$P = Q$. Dengan demikian satu-satunya titik tetap adalah $P$ sendiri, dan
\cref{thm:b3:groups:pfixed} memberikan $n_p = \abs{\mathrm{Syl}_p}
\equiv \abs{\mathrm{Syl}_p^P} = 1 \pmod p$.
\end{proof}

\begin{method}\label{met:b3:groups:sylowmethod}
Untuk menganalisis grup berorde $n = p^a m$: daftarlah pembagi
$m$ yang kongruen dengan $1$ modulo $p$ --- itulah calon-calon bagi
$n_p$. Jika satu-satunya calon adalah $1$, subgrup Sylow-$p$ bersifat
\omterm{def:b3:groups:normal}{normal}. Jika $n_p > 1$ terpaksa bernilai kecil, gunakan \omterm{def:b3:groups:action}{aksi} konjugasi pada
$\mathrm{Syl}_p$ untuk memperoleh morfisma $G \to S_{n_p}$ berkernel kecil.
Lalu hitunglah unsurnya: dua subgrup Sylow-$p$ berbeda yang berorde
\emph{prima} $p$ berpotongan trivial, jadi bersama-sama keduanya membawa
$n_p(p-1)$ unsur yang ordenya tepat $p$; pencacahan yang tumpang tindih untuk
beberapa bilangan prima sering memaksa munculnya kontradiksi
(\cref{exo:b3:groups:7}).
\end{method}

\begin{example}\label{ex:b3:groups:pq}
Misalkan $\abs G = pq$ dengan $p < q$ bilangan prima dan $p \nmid q - 1$. Maka
$n_q \mid p$ dan $n_q \equiv 1 \bmod q$ memaksa $n_q = 1$ (karena $p <
q$); $n_p \mid q$ dan $n_p \equiv 1 \bmod p$ memaksa $n_p = 1$ (karena $q
\not\equiv 1 \bmod p$). Misalkan $P, Q$ dua \omterm{def:b3:groups:sylow}{subgrup Sylow} \omterm{def:b3:groups:normal}{normal} itu: $P
\cap Q = \{e\}$ (ordenya relatif prima), sehingga $\abs{PQ} = pq$
(\cref{exo:b3:groups:4}) dan $G \cong P \times Q \cong \Z/p\Z
\times \Z/q\Z \cong \Z/pq\Z$ menurut
\cref{prop:b3:groups:direct} di bawah. Setiap grup berorde $15$,
$33$, $35$, \dots\ bersifat siklik. Kasus yang dikecualikan, $p \mid q - 1$,
menghasilkan tepat satu grup lagi yang tak abelian --- lihat soal
akhir pekan (\cref{pb:b3:groups:1}).
\end{example}

\begin{example}[Sensus Sylow yang lengkap: $S_4$]
\label{ex:b3:groups:s4sylow}
Mari kita jalankan metode ini pada $G = S_4$, $\abs G = 24 = 2^3\cdot3$.
\emph{Sylow $3$:} $n_3 \mid 8$, $n_3 \equiv 1 \bmod 3$, sehingga $n_3
\in \{1, 4\}$; karena $\langle(123)\rangle$ dan
$\langle(124)\rangle$ berbeda, $n_3 = 4$ --- keempat
subgrup $\langle(abc)\rangle$, satu untuk setiap himpunan bagian beranggota $3$
$\{a, b, c\}$, mencakup ke-$8$ siklus tiga. Menurut Sylow~II
subgrup-subgrup itu saling konjugat, dan morfisma konjugasi $S_4 \to
S_{\mathrm{Syl}_3} \cong S_4$ di sini merupakan isomorfisma (kernelnya
termuat di $N = N_G(\langle(123)\rangle)$ yang berorde $24/4 =
6$, dan \omterm{def:b3:groups:normal}{subgrup normal} dari $S_4$ di dalam grup serupa-$S_3$ ini, yaitu $N$, haruslah
trivial: ia akan terdiri atas permutasi genap yang menetapkan keempat
\omterm{def:b3:groups:sylow}{subgrup Sylow}, dan hanya $e$ yang berbuat demikian). \emph{Sylow $2$:} $n_2 \mid 3$, $n_2
\equiv 1 \bmod 2$: $n_2 \in \{1, 3\}$. Subgrup $D =
\langle(1234), (13)\rangle$ berorde $8$ (sebuah dihedral $D_4$:
simetri persegi dengan titik sudut $1, 2, 3, 4$), tidak
\omterm{def:b3:groups:normal}{normal} ($(12)(1234)(12) = (2134)$ membangun subgrup siklus-$4$
yang lain), sehingga $n_2 = 3$: ketiga salinan $D_4$
bersesuaian dengan tiga cara memasangkan $4$ titik menjadi sebuah
``persegi''. Perhatikan pelajaran dari sensus ini: $\abs{S_4} = 24$
menyisakan ruang bagi kedua Sylow untuk gagal \omterm{def:b3:groups:normal}{normal}, dan keduanya memang gagal ---
bandingkan dengan orde $12$, yang pencacahannya memaksa salah satunya \omterm{def:b3:groups:normal}{normal}
(Bagian~IV dari \cref{pb:b3:groups:1}).
\end{example}

\section{Hasil kali: langsung dan semilangsung}

\begin{proposition}[Mengenali hasil kali langsung]\label{prop:b3:groups:direct}
Misalkan $H, K \trianglelefteq G$ dengan $H \cap K = \{e\}$ dan $HK = G$.
Maka $(h,k) \mapsto hk$ adalah isomorfisma $H \times K \to
G$.\index{hasil kali langsung}
\end{proposition}

\begin{proof}
Untuk $h \in H$, $k \in K$, \omterm{def:b3:groups:derived}{komutator} $hkh^{-1}k^{-1}$ terletak di
$K$ (bacalah sebagai $(hkh^{-1})k^{-1}$, dengan memakai kenormalan $K$) dan di
$H$ (bacalah sebagai $h(kh^{-1}k^{-1})$): jadi ia sama dengan $e$, sehingga $H$ dan $K$
komutatif unsur demi unsur dan pemetaan itu morfisma. Ia surjektif
karena $HK = G$, dan injektif karena $hk = e$ memberikan $h = k^{-1} \in
H \cap K = \{e\}$.
\end{proof}

Kenormalan \emph{kedua} faktornyalah yang paling sering gagal: di $S_3
= \langle (1\,2\,3)\rangle \,\langle(1\,2)\rangle$ kedua faktor
berpotongan trivial dan membangun seluruh grup, namun $S_3 \not\cong \Z/3\Z \times
\Z/2\Z$. Gagasan yang tepat ketika hanya satu faktor yang \omterm{def:b3:groups:normal}{normal} adalah:

\begin{definition}\label{def:b3:groups:semidirect}
Misalkan $H$, $K$ grup dan $\varphi \colon K \to
\operatorname{Aut}(H)$ sebuah morfisma. Adapun \emph{hasil kali
semilangsung}\index{hasil kali semilangsung} $H \rtimes_\varphi K$ adalah
himpunan $H \times K$ yang dilengkapi dengan
\[
(h, k)\,(h', k') = \bigl(h\,\varphi(k)(h'),\; kk'\bigr).
\]
\end{definition}

\begin{proposition}\label{prop:b3:groups:semidirect}
$H \rtimes_\varphi K$ adalah grup; $H \times \{e\}$ \omterm{def:b3:groups:normal}{subgrup normal}
yang isomorfik dengan $H$, dan $\{e\} \times K$ subgrup yang isomorfik
dengan $K$; keduanya berpotongan trivial dan membangun seluruh grup. Sebaliknya, jika $G =
NK$ dengan $N \trianglelefteq G$, $K \leq G$ dan $N \cap K = \{e\}$,
maka $G \cong N \rtimes_\varphi K$ untuk $\varphi(k) = (n \mapsto
knk^{-1})$.
\end{proposition}

\begin{proof}
Pemeriksaan langsung: keasosiatifan menyusut menjadi $\varphi(kk') =
\varphi(k)\circ\varphi(k')$ dan kenyataan bahwa setiap $\varphi(k)$ morfisma;
identitasnya $(e,e)$ dan $(h,k)^{-1} =
\bigl(\varphi(k^{-1})(h^{-1}), k^{-1}\bigr)$. Proyeksi $(h,k)
\mapsto k$ adalah morfisma pada $K$ dengan kernel $H \times \{e\}$,
yang karenanya \omterm{def:b3:groups:normal}{normal}. Sebaliknya: setiap $g \in G$ ditulis
secara \emph{tunggal} sebagai $nk$ dengan $n \in N$, $k \in K$ (keberadaan: $G =
NK$; ketunggalan: $nk = n'k'$ memberikan $n'^{-1}n = k'k^{-1} \in N \cap
K$), dan
\[
(nk)(n'k') = n\,(kn'k^{-1})\;kk'
\]
menunjukkan bahwa $nk \mapsto (n, k)$ mengangkut operasi $G$ menjadi operasi
$N \rtimes_\varphi K$.
\end{proof}

\begin{example}\label{ex:b3:groups:semidirectexamples}
(a) Adapun \emph{grup dihedral}\index{grup dihedral} $D_n$ ($n \geq
3$), yang terdiri atas $2n$ simetri segi-$n$ beraturan: rotasinya membentuk
\omterm{def:b3:groups:normal}{subgrup normal} berindeks $2$, sembarang pencerminan membangun sebuah
komplemen, dan mengonjugasikan sebuah rotasi dengan pencerminan membalik arahnya:
$D_n \cong \Z/n\Z \rtimes_\varphi \Z/2\Z$ dengan $\varphi(1) = (x
\mapsto -x)$.
(b) Grup afin sebuah garis, $\{x \mapsto ax + b : a \in
K^\times,\, b \in K\} \cong K \rtimes K^\times$: translasinya
\omterm{def:b3:groups:normal}{normal}, homotetinya sebuah komplemen.
(c) $S_n \cong A_n \rtimes \Z/2\Z$ (komplemennya: sembarang
transposisi).
(d) \omterm{pb:b3:groups:1}{Grup kuaternion} $Q_8$ \emph{bukan} \omterm{def:b3:groups:semidirect}{hasil kali semilangsung}
dua subgrup sejati: setiap subgrup tak trivial memuat $-1$
(\cref{pb:b3:groups:1}), jadi tidak ada dua subgrup sejati yang berpotongan
trivial.
\end{example}

\section{Grup terselesaikan; kesederhanaan \texorpdfstring{$A_n$}{A_n}}

\begin{definition}\label{def:b3:groups:derived}
Yang disebut \emph{komutator}\index{komutator} dari $x, y \in G$ adalah $[x,y]
= xyx^{-1}y^{-1}$; \emph{subgrup komutator}\index{subgrup
komutator} $D(G)$ adalah subgrup yang dibangun oleh semua komutator.
\emph{Deret komutator} adalah $D^0(G) = G$, $D^{i+1}(G) = D(D^i(G))$,
dan $G$ disebut \emph{terselesaikan}\index{grup terselesaikan} jika $D^n(G) =
\{e\}$ untuk suatu $n$.
\end{definition}

\begin{proposition}\label{prop:b3:groups:derived}
$D(G)$ bersifat \omterm{def:b3:groups:normal}{normal} (bahkan stabil terhadap setiap automorfisma), $G/D(G)$
abelian, dan untuk $N \trianglelefteq G$: $G/N$ abelian $\iff
D(G) \subseteq N$. Lebih lanjut, $G$ \omterm{def:b3:groups:derived}{terselesaikan} jika dan hanya jika ada rantai
$G = G_0 \trianglerighteq G_1 \trianglerighteq \dots
\trianglerighteq G_n = \{e\}$ dengan setiap $G_{i+1} \trianglelefteq
G_i$ dan setiap kuosien $G_i/G_{i+1}$ abelian. Subgrup dan
\omterm{thm:b3:groups:quotient}{grup kuosien} dari \omterm{def:b3:groups:derived}{grup terselesaikan} juga \omterm{def:b3:groups:derived}{terselesaikan}; sebaliknya, jika $N$ dan
$G/N$ \omterm{def:b3:groups:derived}{terselesaikan}, maka $G$ pun \omterm{def:b3:groups:derived}{terselesaikan}.
\end{proposition}

\begin{proof}
Sebuah automorfisma $\alpha$ memetakan $[x,y]$ ke $[\alpha x, \alpha y]$: ia
mempermutasikan \omterm{def:b3:groups:derived}{komutator}, sehingga mempertahankan subgrup yang dibangun olehnya;
konjugasi adalah automorfisma, dari sinilah kenormalannya. Di $G/D(G)$,
$\bar x\bar y \bar x^{-1}\bar y^{-1} = \overline{[x,y]} = \bar e$:
\omterm{thm:b3:groups:quotient}{grup kuosiennya} abelian. Jika $G/N$ abelian maka setiap $[x,y] \in
N$, sehingga $D(G) \subseteq N$; sebaliknya jika $D(G) \subseteq N$ maka
$G/N$, yang merupakan \omterm{thm:b3:groups:quotient}{grup kuosien} dari $G/D(G)$ yang abelian menurut teorema
isomorfisma ketiga, juga abelian.

Jika $G$ \omterm{def:b3:groups:derived}{terselesaikan}, deret \omterm{def:b3:groups:derived}{komutatornya} adalah rantai semacam itu. Sebaliknya,
diberikan sebuah rantai, $D^i(G) \subseteq G_i$ dengan induksi: $G_i/G_{i+1}$
abelian memberikan $D(G_i) \subseteq G_{i+1}$, jadi $D^{i+1}(G) = D(D^i G)
\subseteq D(G_i) \subseteq G_{i+1}$; akibatnya $D^n(G) = \{e\}$.

Sifat menurun: $D^i(H) \subseteq D^i(G)$ untuk $H \leq G$ (induksi),
dan $D^i(G/N) = \pi(D^i(G))$ sebab $\pi$ memetakan \omterm{def:b3:groups:derived}{komutator} pada
\omterm{def:b3:groups:derived}{komutator}; dari sini diperoleh pernyataan untuk subgrup dan \omterm{thm:b3:groups:quotient}{grup kuosien}.
Perluasan: jika $D^m(G/N) = \{e\}$ maka $D^m(G) \subseteq N$, dan
$D^n(N) = \{e\}$ memberikan $D^{m+n}(G) = D^n(D^m(G)) \subseteq D^n(N)
= \{e\}$.
\end{proof}

\begin{example}\label{ex:b3:groups:solvableexamples}
Grup abelian \omterm{def:b3:groups:derived}{terselesaikan}. Grup-$p$ \omterm{def:b3:groups:derived}{terselesaikan}, dengan induksi
pada ordenya: $Z(G) \neq \{e\}$ dan $G/Z(G)$ adalah grup-$p$ yang
lebih kecil. $S_3$ dan $S_4$ \omterm{def:b3:groups:derived}{terselesaikan}: $S_4 \trianglerighteq A_4
\trianglerighteq V \trianglerighteq \{e\}$, dengan $V = \{e,
(1\,2)(3\,4), (1\,3)(2\,4), (1\,4)(2\,3)\}$ adalah grup Klein yang terdiri atas
transposisi ganda (\omterm{def:b3:groups:normal}{normal} di $S_4$: ia gabungan \omterm{ex:b3:groups:actions}{kelas
konjugasi}), dan kuosien-kuosiennya abelian: $\Z/2\Z$, $\Z/3\Z$, $V$. Pada
\cref{ch:b3:galois}, ``persamaan umum berderajat $n$
\omterm{def:b3:groups:derived}{terselesaikan} dengan radikal'' akan benar-benar \emph{berarti} ``$S_n$ adalah
\omterm{def:b3:groups:derived}{grup terselesaikan}''. Dari sinilah pentingnya definisi berikut.
\end{example}

\begin{definition}\label{def:b3:groups:simple}
Grup $G \neq \{e\}$ disebut \emph{sederhana}\index{grup sederhana} jika
\omterm{def:b3:groups:normal}{subgrup normalnya} hanyalah $\{e\}$ dan $G$. Grup sederhana yang tak
abelian tidak \omterm{def:b3:groups:derived}{terselesaikan}: $D(G) \trianglelefteq G$ tidak sama dengan $\{e\}$
(sebab kalau tidak $G$ abelian), jadi $D(G) = G$ dan deret \omterm{def:b3:groups:derived}{komutatornya}
konstan. Grup sederhana yang abelian tepat berupa $\Z/p\Z$ dengan $p$
prima (grup abelian sederhana jika dan hanya jika ia tak punya subgrup sejati yang tak
trivial, jika dan hanya jika ia siklik berorde prima menurut Lagrange).
\end{definition}

\begin{lemma}\label{lem:b3:groups:threecycles}
Untuk $n \geq 3$, $A_n$ dibangun oleh siklus-$3$; untuk $n \geq 5$,
semua siklus-$3$ saling konjugat \emph{di dalam $A_n$}.
\end{lemma}

\begin{proof}
Unsur $A_n$ adalah hasil kali sejumlah genap
transposisi; pasangkanlah transposisi itu lalu gunakan (komposisi dari kanan ke kiri)
\[
(a\,b)(c\,d) = (a\,c\,b)(a\,c\,d), \qquad
(a\,b)(b\,c) = (a\,b\,c), \qquad
(a\,b)(a\,b) = e
\]
berturut-turut untuk pasangan yang saling lepas, yang bertindih dan yang sama: setiap pasangan
transposisi adalah hasil kali siklus-$3$.

Kekonjugatan: $\sigma(a\,b\,c)\sigma^{-1} = (\sigma a\; \sigma b\;
\sigma c)$, jadi dua siklus-$3$ mana pun saling konjugat oleh suatu $\sigma \in
S_n$. Jika $\sigma$ ganjil, gantilah dengan $\sigma' = \sigma (d\,e)$
dengan $d, e$ dua titik di luar $\{a, b, c\}$ --- keduanya ada
karena $n \geq 5$; maka $\sigma'$ genap dan $\sigma'(a\,b\,c)
\sigma'^{-1} = \sigma(a\,b\,c)\sigma^{-1}$, sebab $(d\,e)$ komutatif
dengan $(a\,b\,c)$.
\end{proof}

\begin{theorem}[Kesederhanaan grup alternasi]\label{thm:b3:groups:ansimple}
$A_n$ \omterm{def:b3:groups:simple}{sederhana} untuk $n \geq 5$.\index{grup alternasi}
\end{theorem}

\begin{proof}
Misalkan $N \trianglelefteq A_n$, $N \neq \{e\}$. Menurut
\cref{lem:b3:groups:threecycles} cukup ditunjukkan bahwa $N$
memuat \emph{satu} siklus-$3$: kenormalan dan kekonjugatan
siklus-$3$ di dalam $A_n$ lalu menempatkan semua siklus-$3$ di $N$, jadi $N = A_n$.

Untuk $\rho \in S_n$ tulis $F(\rho) = \{x : \rho(x) \neq x\}$ sebagai
\emph{pendukungnya} dan $f(\rho) = \abs{F(\rho)}$. Pilihlah $\sigma \in N
\setminus \{e\}$ dengan $f(\sigma)$ \emph{minimal}. Perhatikan bahwa
permutasi genap yang tak trivial mempunyai $f \geq 3$, dan bahwa $f(\sigma) =
4$ mustahil bagi $\sigma \in A_n$ kecuali jika $\sigma$ transposisi
ganda (siklus-$4$ bersifat ganjil). Kita tunjukkan bahwa $\sigma$ adalah
siklus-$3$.

\emph{Kasus A: $\sigma$ hasil kali transposisi yang saling lepas},
katakanlah $\sigma = (a\,b)(c\,d)\cdots$ dengan $f(\sigma) \geq 4$. Ambil $e'
\notin \{a, b, c, d\}$ (bisa dilakukan karena $n \geq 5$), tetapkan $\tau =
(c\,d\,e')$ dan
\[
\sigma' = \tau\sigma\tau^{-1}\,\sigma^{-1} \in N
\qquad (\tau\sigma\tau^{-1} \in N \text{ menurut kenormalan}).
\]
Karena $\sigma\tau^{-1}\sigma^{-1} = (\sigma c\;\sigma e'\;\sigma
d) = (d\;\sigma e'\;c)$ (memakai $\sigma c = d$, $\sigma d = c$),
diperoleh $\sigma' = (c\,d\,e')(d\;\sigma e'\;c)$.

Jika $\sigma e' = e'$ (yang khususnya berlaku ketika $f(\sigma) =
4$, yaitu $\sigma = (a\,b)(c\,d)$): maka $(d\,e'\,c) =
(c\,d\,e')$ dan $\sigma' = (c\,d\,e')^2 = (c\,e'\,d)$, sebuah
siklus-$3$ yang terletak di $N$, dengan $f(\sigma') = 3 < 4 \leq f(\sigma)$
--- bertentangan dengan keminimalan.

Jika $\sigma e' \neq e'$: maka $\sigma e' \notin \{a, b, c, d,
e'\}$ (sebab $\sigma$ menukar $a,b$ dan $c,d$, sedangkan $e' \notin \{a,b,c,d\}$
dan $\sigma$ injektif), jadi $\sigma$ menggerakkan keenam titik $a, b,
c, d, e', \sigma e'$: $f(\sigma) \geq 6$. Di pihak lain
$\sigma'$, hasil kali dua siklus-$3$ yang pendukungnya di dalam $\{c, d,
e', \sigma e'\}$, memenuhi $f(\sigma') \leq 4$; dan $\sigma' \neq
e$, sebab $\sigma'(d) = \tau\sigma\tau^{-1}(c) = \tau\sigma(e') =
\sigma e' \neq d$ ($\tau$ menetapkan $\sigma e' \notin \{c,d,e'\}$).
Jadi $\sigma' \in N \setminus\{e\}$ dengan $f(\sigma') \leq 4 <
f(\sigma)$: keminimalan dilanggar.

\emph{Kasus B: suatu siklus dari $\sigma$ berpanjang $\geq 3$}, katakanlah
$\sigma(a) = b$, $\sigma(b) = c$ dengan $a, b, c$ berbeda. Jika
$\sigma$ tepat siklus-$3$ ini, pembuktian selesai. Jika tidak,
$f(\sigma) \geq 5$ (kasus $f(\sigma) = 4$ dengan siklus berpanjang $\geq3$
berupa siklus-$4$ yang ganjil, dan sudah dikecualikan), sehingga kita boleh memilih $d, e' \in
F(\sigma) \setminus \{a, b, c\}$. Tetapkan $\tau = (c\,d\,e')$ dan
$\sigma' = \tau\sigma\tau^{-1}\sigma^{-1} \in N$. Seperti sebelumnya
$\sigma' = (c\,d\,e')\,(\sigma c\;\sigma e'\;\sigma d)$ hanya menggerakkan
titik-titik di
\[
M = \{c, d, e'\} \cup \{\sigma c, \sigma d, \sigma e'\}
\subseteq F(\sigma)
\]
(peta dari titik yang bergerak ikut bergerak: $\sigma(x) \ne x$ mengakibatkan
$\sigma(\sigma x) \neq \sigma x$, karena $\sigma$ injektif). Selain itu
$b \notin M$: kelima titik $a, b, c, d, e'$ berbeda, jadi $b
\notin \{c, d, e'\}$; sedangkan $b \in \{\sigma c, \sigma d, \sigma
e'\}$ akan memaksa $a \in \{c, d, e'\}$ (terapkan $\sigma^{-1}$, dengan memakai
$\sigma a = b$), dan itu salah. Karena itu $\sigma'$ menetapkan $b$, padahal
$\sigma$ menggerakkan $b$; lagi pula
$F(\sigma') \subseteq F(\sigma)$. Akhirnya $\sigma' \neq e$:
$\sigma^{-1}(c) = b$, $\tau^{-1}(b) = b$, $\sigma(b) = c$,
$\tau(c) = d$, sehingga $\sigma'(c) = d \neq c$. Jadi $\sigma' \in N
\setminus \{e\}$ dengan $f(\sigma') \leq f(\sigma) - 1$,
dan keminimalan kembali dilanggar.

Karena kedua kasus mustahil, $\sigma$ pastilah siklus-$3$.
\end{proof}

\begin{corollary}\label{cor:b3:groups:snnotsolvable}
Untuk $n \geq 5$: $A_n$ dan $S_n$ tidak \omterm{def:b3:groups:derived}{terselesaikan}, dan satu-satunya
\omterm{def:b3:groups:normal}{subgrup normal} $S_n$ adalah $\{e\}$, $A_n$ dan $S_n$.
\end{corollary}

\begin{proof}
$A_n$ \omterm{def:b3:groups:simple}{sederhana} dan tak abelian, sehingga tidak \omterm{def:b3:groups:derived}{terselesaikan}
(\cref{def:b3:groups:simple}); grup yang memuat subgrup tak
\omterm{def:b3:groups:derived}{terselesaikan} juga tidak \omterm{def:b3:groups:derived}{terselesaikan} (\cref{prop:b3:groups:derived}). Misalkan $N
\trianglelefteq S_n$: maka $N \cap A_n \trianglelefteq A_n$ sama dengan
$\{e\}$ atau $A_n$. Jika $N \cap A_n = A_n$, maka $A_n \subseteq N$
dan $N \in \{A_n, S_n\}$ karena indeksnya. Jika $N \cap A_n = \{e\}$,
pembatasan pada $N$ dari proyeksi $S_n \to S_n/A_n \cong
\Z/2\Z$ bersifat injektif, jadi $\abs N \leq 2$; jika $N = \{e, \sigma\}$,
kenormalan membuat \omterm{ex:b3:groups:actions}{kelas konjugasi} $\sigma$ sama dengan
$\{\sigma\}$, yakni $\sigma \in Z(S_n)$. Padahal $Z(S_n) = \{e\}$ untuk
$n \geq 3$: jika $\sigma \ne e$ memindahkan $a$ ke $b \neq a$, pilih $c
\notin \{a, b\}$; maka $(b\,c)\sigma(b\,c)^{-1}$ mengirim $a$ ke $c
\ne b$, jadi ia berbeda dari $\sigma$. Karena itu $N = \{e\}$.
\end{proof}

\begin{theorem}[Jordan--Hölder]\label{thm:b3:groups:jordanholder}
Setiap grup berhingga $G \neq \{e\}$ mempunyai \emph{deret
komposisi}\index{deret komposisi}
\[
\{e\} = G_0 \trianglelefteq G_1 \trianglelefteq \cdots
\trianglelefteq G_r = G,
\qquad G_{i}/G_{i-1} \text{ sederhana},
\]
dan multihimpunan \emph{faktor komposisi} $G_i/G_{i-1}$, sampai
isomorfisma, tidak bergantung pada deret yang dipilih. Grup
berhingga \omterm{def:b3:groups:derived}{terselesaikan} jika dan hanya jika semua faktor komposisinya siklik
berorde prima.\index{teorema Jordan--Hölder}
\end{theorem}

\begin{proof}
\emph{Keberadaan}: induksi pada $\abs G$. Jika $G$ \omterm{def:b3:groups:simple}{sederhana}, ambil
$\{e\} \trianglelefteq G$. Jika tidak, pilihlah \omterm{def:b3:groups:normal}{subgrup normal} sejati yang
maksimal, sebutlah $N$ (subgrupnya berhingga banyaknya); $G/N$ \omterm{def:b3:groups:simple}{sederhana}
menurut teorema korespondensi (\omterm{def:b3:groups:normal}{subgrup normal} sejati yang tak trivial
dari $G/N$ akan terangkat menjadi \omterm{def:b3:groups:normal}{subgrup normal} $G$ yang terletak tepat di antara
$N$ dan $G$). Tambahkan $N \trianglelefteq G$ pada \omterm{thm:b3:groups:jordanholder}{deret komposisi}
milik $N$.

\emph{Ketunggalan}: induksi pada $\abs G$, dengan kasus $G$ \omterm{def:b3:groups:simple}{sederhana}
sudah jelas. Ambil dua \omterm{thm:b3:groups:jordanholder}{deret komposisi}, dengan suku kedua dari akhir
$M \trianglelefteq G$ dan $N \trianglelefteq G$ (jadi $G/M$, $G/N$
\omterm{def:b3:groups:simple}{sederhana}). Jika $M = N$, simpulkan dengan induksi yang diterapkan pada $M$.
Jika tidak, $MN$, yang \omterm{def:b3:groups:normal}{normal} di $G$ dan memuat $M$ secara sejati, sama dengan
$G$ (sebab $M$ maksimal \omterm{def:b3:groups:normal}{normal}: \omterm{def:b3:groups:normal}{subgrup normal} apa pun dengan $M \subsetneq L
\subsetneq G$ akan terpetakan menjadi \omterm{def:b3:groups:normal}{subgrup normal} sejati yang tak trivial dari
$G/M$ yang \omterm{def:b3:groups:simple}{sederhana}). Teorema isomorfisma kedua memberikan
\[
G/M = MN/M \cong N/(M \cap N),
\qquad
G/N = MN/N \cong M/(M \cap N).
\]
Tetapkan $K = M \cap N$ ($\trianglelefteq G$) lalu pilih sebuah \omterm{thm:b3:groups:jordanholder}{deret
komposisi} bagi $K$. Maka $M$ membawa dua \omterm{thm:b3:groups:jordanholder}{deret komposisi}: deret
aslinya, dan deret $K$ yang disusul $K \trianglelefteq
M$ (kuosiennya $M/K \cong G/N$ \omterm{def:b3:groups:simple}{sederhana}). Dengan induksi (yang diterapkan
pada $M$), faktor deret asli $M$ adalah
$\{\text{faktor } K\} \cup \{G/N\}$; demikian pula untuk $N$. Karena itu
kedua deret $G$ mempunyai faktor
\[
\{\text{faktor } K\} \;\cup\; \{\,G/N,\; G/M\,\},
\]
yaitu multihimpunan yang sama.

\omterm{def:b3:groups:derived}{Keterselesaian}: jika semua faktornya $\Z/p_i\Z$, deret itu adalah rantai
dengan kuosien abelian, jadi $G$ \omterm{def:b3:groups:derived}{terselesaikan}
(\cref{prop:b3:groups:derived}). Sebaliknya, faktor komposisi
sebuah \omterm{def:b3:groups:derived}{grup terselesaikan} juga \omterm{def:b3:groups:derived}{terselesaikan} (ia \omterm{thm:b3:groups:quotient}{grup kuosien} dari sebuah subgrup) sekaligus
\omterm{def:b3:groups:simple}{sederhana}; \omterm{def:b3:groups:simple}{grup sederhana} yang \omterm{def:b3:groups:derived}{terselesaikan} pastilah abelian ($D(G) \ne G$ memaksa
$D(G) = \{e\}$), jadi berupa $\Z/p\Z$.
\end{proof}

\begin{remark}\label{rem:b3:groups:jhmeaning}
Jordan--Hölder mengatakan bahwa setiap grup berhingga dibangun dari \omterm{def:b3:groups:simple}{grup
sederhana}, dengan daftar komponen yang terdefinisi dengan baik --- sebuah aritmetika grup
yang di dalamnya \omterm{def:b3:groups:simple}{grup sederhana} berperan sebagai bilangan prima, dan yang di dalamnya \emph{cara
komponen itu direkatkan} (data perluasan, seperti pada \omterm{def:b3:groups:semidirect}{hasil kali semilangsung})
menggantikan perkalian biasa. Klasifikasi \omterm{def:b3:groups:simple}{grup sederhana}
berhingga --- yang siklik $\Z/p\Z$, yang alternasi $A_{n \geq
5}$, enam belas keluarga bertipe Lie, dan $26$ grup sporadis --- adalah
salah satu monumen matematika abad kedua puluh; pembuktiannya,
yang terbentang di sekitar sepuluh ribu halaman jurnal, jauh melampaui
kuliah ini.
\end{remark}

\begin{omfigure}
\begin{tikzpicture}[xscale=1.55, yscale=1.4]
  \tikzset{sg/.style={draw, rounded corners=2pt, inner sep=3pt,
      fill=omDef!6, font=\small}}
  \node[sg, fill=omThm!10] (G)   at (0, 3)    {$D_4$};
  \node[sg] (K1)  at (-1.6, 2)   {$\langle r^2, s\rangle$};
  \node[sg] (R)   at (0, 2)      {$\langle r\rangle$};
  \node[sg] (K2)  at (1.6, 2)    {$\langle r^2, rs\rangle$};
  \node[sg] (S1)  at (-2.4, 1)   {$\langle s\rangle$};
  \node[sg] (S2)  at (-1.2, 1)   {$\langle r^2 s\rangle$};
  \node[sg, fill=omProp!12] (Z) at (0, 1) {$\langle r^2\rangle$};
  \node[sg] (S3)  at (1.2, 1)    {$\langle rs\rangle$};
  \node[sg] (S4)  at (2.4, 1)    {$\langle r^3 s\rangle$};
  \node[sg] (E)   at (0, 0)      {$\{e\}$};
  \draw (G) -- (K1); \draw (G) -- (R); \draw (G) -- (K2);
  \draw (K1) -- (S1); \draw (K1) -- (S2); \draw (K1) -- (Z);
  \draw (R) -- (Z);
  \draw (K2) -- (Z); \draw (K2) -- (S3); \draw (K2) -- (S4);
  \draw (S1) -- (E); \draw (S2) -- (E); \draw (Z) -- (E);
  \draw (S3) -- (E); \draw (S4) -- (E);
\end{tikzpicture}

{\small Kesepuluh subgrup dari \omterm{ex:b3:groups:semidirectexamples}{grup dihedral} $D_4 = \langle r,
s \mid r^4 = s^2 = e,\ srs^{-1} = r^{-1}\rangle$. Ketiga
subgrup yang berindeks $2$ (baris tengah) bersifat \omterm{def:b3:groups:normal}{normal}, demikian pula pusatnya
$\langle r^2\rangle$ (yang disorot); keempat subgrup pencerminan
terbagi atas dua \omterm{ex:b3:groups:actions}{kelas konjugasi} beranggota dua. Rantai dari bawah ke atas
memberikan \omterm{thm:b3:groups:jordanholder}{deret komposisi}, misalnya $\{e\} \trianglelefteq \langle
r^2\rangle \trianglelefteq \langle r\rangle \trianglelefteq D_4$:
faktornya $\Z/2\Z, \Z/2\Z, \Z/2\Z$ --- selalu multihimpunan yang sama, seperti
yang dituntut Jordan--Hölder.}
\end{omfigure}

\section{Latihan}

\begin{exercise}[$\star$]\label{exo:b3:groups:1}
(a) Tunjukkan bahwa setiap subgrup berindeks $2$ bersifat \omterm{def:b3:groups:normal}{normal}.
(b) Tunjukkan bahwa jika $[G : H] = 2$, maka $x^2 \in H$ untuk setiap $x \in
G$.
(c) Simpulkan bahwa $A_4$ tidak mempunyai subgrup berorde $6$: konvers
teorema Lagrange gagal. \emph{(Hitunglah kuadrat siklus-$3$.)}
\end{exercise}

\begin{exercise}[$\star$]\label{exo:b3:groups:2}
Tunjukkan bahwa jika $G/Z(G)$ siklik maka $G$ abelian. Simpulkan lagi
bahwa setiap grup berorde $p^2$ abelian, lalu tampilkan, untuk setiap
bilangan prima $p$, sebuah grup tak abelian berorde $p^3$. \emph{(Pikirkan
matriks segitiga atas dengan diagonal satuan atas $\mathbb F_p$.)}
\end{exercise}

\begin{exercise}[$\star$]\label{exo:b3:groups:3}
(a) Tunjukkan bahwa $\operatorname{Aut}(\Z/n\Z) \cong (\Z/n\Z)^\times$.
(b) Tunjukkan bahwa automorfisma dalam $\iota_g \colon x \mapsto
gxg^{-1}$ membentuk \omterm{def:b3:groups:normal}{subgrup normal} $\operatorname{Inn}(G)
\trianglelefteq \operatorname{Aut}(G)$, dengan
$\operatorname{Inn}(G) \cong G/Z(G)$.
\end{exercise}

\begin{exercise}[$\star\star$]\label{exo:b3:groups:4}
Misalkan $H, K$ subgrup sebuah grup berhingga $G$.
(a) Buktikan \emph{rumus hasil kali} $\abs{HK}\,\abs{H\cap K} =
\abs H\, \abs K$, dengan mencacah serat pemetaan $H \times K
\to HK$, $(h,k) \mapsto hk$.
(b) Tunjukkan bahwa $HK$ subgrup jika dan hanya jika $HK = KH$ (yang otomatis jika salah
satu dari keduanya \omterm{def:b3:groups:normal}{normal}).
(c) Jika $H, K \trianglelefteq G$ dan $H \cap K = \{e\}$, tunjukkan bahwa
$hk = kh$ untuk setiap $h \in H$, $k \in K$.
\end{exercise}

\begin{exercise}[$\star\star$]\label{exo:b3:groups:5}
(Lema pencacahan Burnside) Sebuah grup berhingga $G$ beraksi pada himpunan
berhingga $X$. Tunjukkan bahwa banyaknya \omterm{def:b3:groups:action}{orbit} sama dengan \emph{rata-rata
banyak titik tetap}:
\[
\#\{\text{orbit}\} = \frac{1}{\abs G}\sum_{g \in G}
\abs{\operatorname{Fix}(g)},
\qquad \operatorname{Fix}(g) = \{x \in X : g \cdot x = x\},
\]
dengan mencacah himpunan $\{(g,x) : g\cdot x = x\}$ dengan dua cara.
Penerapan: $\Z/p\Z$ ($p$ prima) beraksi melalui rotasi pada kalung berisi
$p$ manik dengan $a$ warna yang tersedia; simpulkan teorema kecil
Fermat $a^p \equiv a \pmod p$.
\end{exercise}

\begin{exercise}[$\star$]\label{exo:b3:groups:6}
Dengan memakai teorema Sylow, tunjukkan bahwa setiap grup berorde $15$
siklik dan bahwa setiap grup berorde $45$ abelian.
\end{exercise}

\begin{exercise}[$\star\star$]\label{exo:b3:groups:7}
Tunjukkan bahwa tidak ada grup berorde $30$, dan tidak ada pula yang berorde $56$, yang
\omterm{def:b3:groups:simple}{sederhana}. \emph{(Untuk $30$: jika $n_3 \neq 1$ dan $n_5 \neq 1$, hitunglah
unsur berorde $3$ dan $5$. Untuk $56$: hitunglah unsur
berorde $7$.)}
\end{exercise}

\begin{exercise}[$\star\star$]\label{exo:b3:groups:8}
(a) Misalkan $H \leq G$ berindeks $n$. Tunjukkan bahwa \omterm{def:b3:groups:action}{aksi} $G$ pada
$G/H$ melahirkan morfisma $G \to S_n$ yang kernelnya $\bigcap_{g \in
G} gHg^{-1}$ adalah \omterm{def:b3:groups:normal}{subgrup normal} terbesar dari $G$ yang termuat di
$H$.
(b) Simpulkan: jika $G$ berhingga dan $p$ pembagi prima \emph{terkecil}
dari $\abs G$, maka setiap subgrup berindeks $p$ bersifat \omterm{def:b3:groups:normal}{normal}.
(c) Tunjukkan bahwa setiap \omterm{def:b3:groups:action}{aksi} transitif $G$ pada himpunan $X$
isomorfik dengan \omterm{def:b3:groups:action}{aksi} pada sebuah ruang koset: ada bijeksi $X
\to G/G_x$ yang komutatif dengan kedua \omterm{def:b3:groups:action}{aksi}.
\end{exercise}

\begin{exercise}[$\star\star$]\label{exo:b3:groups:9}
(a) Tunjukkan bahwa $D(G)$ adalah \omterm{def:b3:groups:normal}{subgrup normal} terkecil dari $G$ yang
kuosiennya abelian, dan bahwa setiap morfisma dari $G$ ke sebuah grup
abelian terfaktorkan secara tunggal melalui \emph{abelianisasi}
$G^{\mathrm{ab}} = G/D(G)$.
(b) Hitunglah $D(S_n)$ dan $S_n^{\mathrm{ab}}$ untuk $n \geq 2$, serta
$D(Q_8)$ dan $Q_8^{\mathrm{ab}}$.
\end{exercise}

\begin{exercise}[$\star\star$]\label{exo:b3:groups:10}
Misalkan $G$ sebuah grup-$p$ dan $H \subsetneq G$ subgrup sejati.
Tunjukkan bahwa $H \subsetneq N_G(H)$ (``\omterm{ex:b3:groups:actions}{normalisator} tumbuh''), lalu simpulkan
bahwa setiap subgrup maksimal dari grup-$p$ bersifat \omterm{def:b3:groups:normal}{normal} dan berindeks $p$.
\emph{(Induksi pada $\abs G$, dengan memakai $Z(G) \neq \{e\}$: tinjaulah
secara terpisah $Z(G) \subseteq H$ dan $Z(G) \not\subseteq H$.)}
\end{exercise}

\begin{exercise}[$\star\star\star$]\label{exo:b3:groups:11}
(Kesederhanaan $A_5$, secara langsung)
(a) Tunjukkan bahwa \omterm{ex:b3:groups:actions}{kelas konjugasi} $A_5$ berkardinalitas
$1$, $15$, $20$, $12$, $12$. Perhatikan baik-baik terpecahnya
kelas-$S_5$ dari siklus-$5$: untuk sebuah siklus-$5$ $\sigma$, bandingkan
\omterm{ex:b3:groups:actions}{sentralisator} $\sigma$ di $S_5$ dan di $A_5$.
(b) Simpulkan bahwa $A_5$ \omterm{def:b3:groups:simple}{sederhana}: \omterm{def:b3:groups:normal}{subgrup normal} adalah gabungan
\omterm{ex:b3:groups:actions}{kelas konjugasi}, memuat $e$, dan kardinalitasnya membagi
$60$.
(c) Tunjukkan bahwa \omterm{def:b3:groups:simple}{grup sederhana} berorde $60$ pastilah mempunyai $n_5 =
6$.
\end{exercise}

\begin{exercise}[$\star\star$]\label{exo:b3:groups:12}
(\omterm{ex:b3:groups:actions}{Normalisator} \omterm{def:b3:groups:sylow}{subgrup Sylow} menormalkan dirinya sendiri) Misalkan $P$
subgrup Sylow-$p$ dari grup berhingga $G$ dan $H =
N_G(P)$.
(a) Tunjukkan bahwa $P$ adalah subgrup Sylow-$p$ \emph{tunggal} dari
$H$.
(b) Simpulkan $N_G(H) = H$. \emph{(Untuk $g \in N_G(H)$: $gPg^{-1}$
adalah subgrup Sylow-$p$ dari $H$, jadi $gPg^{-1} = P$.)}
(c) Simpulkan bahwa tidak ada \omterm{ex:b3:groups:actions}{normalisator} Sylow yang termuat di \omterm{def:b3:groups:normal}{subgrup
normal} sejati dari $G$, dan bahwa subgrup maksimal yang memuat
$N_G(P)$ menormalkan dirinya sendiri.
\end{exercise}

\section{Soal: grup berorde tak lebih dari 15}

\begin{problem}[{Soal akhir pekan --- klasifikasi grup
kecil}]\label{pb:b3:groups:1}
Tujuannya adalah klasifikasi lengkap, beserta bukti penuh, atas
grup berorde $\leq 15$ sampai isomorfisma. Orde $1, 2, 3, 5,
7, 11, 13$ sudah tuntas menurut Lagrange (siklik), dan orde $4$ serta
$9$ menurut \cref{thm:b3:groups:pfixed} ditambah analisis di bawah untuk
$p^2$: tersisa $6, 8, 10, 12, 14, 15$.

\medskip
\noindent\textbf{Bagian I --- Perkakas.}
\begin{enumerate}
  \item Tunjukkan bahwa grup yang setiap unsurnya memenuhi $x^2 =
        e$ bersifat abelian; simpulkan bahwa grup berhingga semacam itu berorde
        $2^k$ dan isomorfik dengan $(\Z/2\Z)^k$. \emph{(Pandanglah ia sebagai
        ruang vektor atas $\mathbb F_2$.)}
  \item Tunjukkan bahwa grup berorde $p^2$ isomorfik dengan
        $\Z/p^2\Z$ atau $(\Z/p\Z)^2$. Daftarlah grup abelian berorde
        $8$ sampai isomorfisma: $\Z/8\Z$, $\Z/4\Z \times
        \Z/2\Z$, $(\Z/2\Z)^3$ --- buktikan bahwa daftar itu lengkap dan
        tanpa pengulangan \emph{tanpa} memakai teorema struktur pada
        \cref{ch:b3:modules} (uraikan menurut orde terbesar sebuah
        unsur).
  \item Misalkan $\varphi, \varphi' \colon K \to \operatorname{Aut}(H)$
        dua \omterm{def:b3:groups:action}{aksi}. Tunjukkan bahwa jika $\varphi' = \varphi \circ
        \alpha$ dengan $\alpha \in \operatorname{Aut}(K)$, maka $H
        \rtimes_{\varphi} K \cong H \rtimes_{\varphi'} K$.
  \item Tentukan $\operatorname{Aut}(\Z/n\Z)$ untuk $n = 3, 4, 5,
        7$ secara eksplisit, dan tunjukkan $\operatorname{Aut}\bigl((\Z/2\Z)^2
        \bigr) \cong S_3$.
\end{enumerate}

\medskip
\noindent\textbf{Bagian II --- Orde $2p$ ($6$, $10$, $14$) dan
$pq$.}
\begin{enumerate}[resume]
  \item Misalkan $\abs G = 2p$ dengan $p$ bilangan prima ganjil. Tunjukkan bahwa $G$ mempunyai
        \omterm{def:b3:groups:normal}{subgrup normal} $N = \langle r \rangle$ berorde $p$ dan
        sebuah unsur $s$ berorde $2$ di luar $N$.
  \item Simpulkan $G \cong \Z/p\Z \rtimes_\varphi \Z/2\Z$, dengan
        $\varphi(1) \in \operatorname{Aut}(\Z/p\Z)$ sebuah
        involusi, lalu simpulkan: $G \cong \Z/2p\Z$ atau $G \cong
        D_p$; periksa bahwa keduanya tidak isomorfik. Ini menuntaskan
        orde $6$, $10$, $14$.
  \item Secara lebih umum, misalkan $\abs G = pq$ dengan $p < q$ bilangan prima.
        Tunjukkan bahwa jika $p \nmid q-1$ maka $G$ siklik
        (\cref{ex:b3:groups:pq}), dan bahwa jika $p \mid q - 1$ maka,
        selain $\Z/pq\Z$, ada \emph{tepat satu} grup tak abelian
        $\Z/q\Z \rtimes \Z/p\Z$ sampai isomorfisma --- pakailah pertanyaan
        3 dan kenyataan bahwa $\operatorname{Aut}(\Z/q\Z) \cong
        (\Z/q\Z)^\times$ siklik berorde $q - 1$, yang di sini diterima tanpa bukti
        dan dibuktikan pada \cref{ch:b3:galois} (kesiklikan
        $\mathbb F_q^\times$). Simpulkan untuk orde $15$.
\end{enumerate}

\medskip
\noindent\textbf{Bagian III --- Orde $8$.} Misalkan $G$ tak abelian dan
berorde $8$.
\begin{enumerate}[resume]
  \item Tunjukkan bahwa $G$ mempunyai unsur $r$ berorde $4$ (pakailah
        pertanyaan 1) dan bahwa $N = \langle r\rangle$ \omterm{def:b3:groups:normal}{normal}.
  \item Misalkan $s \notin N$. Tunjukkan bahwa $s^2 \in N$
        (\cref{exo:b3:groups:1}(b)), bahwa $srs^{-1} = r^{-1}$
        (periksalah peta $r$ yang mungkin di bawah konjugasi,
        yang harus berorde 4, lalu kecualikan $srs^{-1} = r$), dan
        bahwa $s^2 \in \{e, r^2\}$ \emph{(apa yang terjadi jika $s^2 = r$
        atau $r^3$? dan mengapa $s^2$ harus komutatif dengan $s$?)}.
  \item Dalam kasus $s^2 = e$, tunjukkan $G \cong D_4$.
  \item Dalam kasus $s^2 = r^2$, tunjukkan bahwa tabel
        perkaliannya sepenuhnya tertentu; grup yang dihasilkan adalah
        \emph{grup kuaternion}\index{grup kuaternion} $Q_8 =
        \{\pm 1, \pm \mathrm i, \pm \mathrm j, \pm \mathrm k\}$,
        $\mathrm i^2 = \mathrm j^2 = \mathrm k^2 = \mathrm i\,
        \mathrm j\,\mathrm k = -1$ (ambil $r = \mathrm i$, $s =
        \mathrm j$). Pastikan bahwa $Q_8$ memang ada, misalnya di dalam
        $GL_2(\C)$ melalui
        \[
        \mathrm i \mapsto \begin{pmatrix} \iu & 0\\ 0 & -\iu
        \end{pmatrix},
        \qquad
        \mathrm j \mapsto \begin{pmatrix} 0 & 1\\ -1 & 0
        \end{pmatrix}.
        \]
  \item Tunjukkan bahwa setiap subgrup tak trivial dari $Q_8$ memuat
        $-1$; simpulkan bahwa setiap subgrup $Q_8$ \omterm{def:b3:groups:normal}{normal}, bahwa
        $D_4 \not\cong Q_8$ (hitunglah unsur berorde $2$), dan
        bahwa $Q_8$ bukan \omterm{def:b3:groups:semidirect}{hasil kali semilangsung} dua subgrup
        sejati.
\end{enumerate}

\medskip
\noindent\textbf{Bagian IV --- Orde $12$.} Misalkan $\abs G = 12$, $P_3
\in \mathrm{Syl}_3(G)$, $P_2 \in \mathrm{Syl}_2(G)$.
\begin{enumerate}[resume]
  \item Tunjukkan $n_3 \in \{1, 4\}$, $n_2 \in \{1, 3\}$, dan bahwa $n_3
        = 4$ memaksa $n_2 = 1$ \emph{(hitunglah unsur berorde
        $3$)}.
  \item Andaikan $n_3 = 4$. \omterm{def:b3:groups:action}{Aksi} konjugasi pada
        $\mathrm{Syl}_3$ memberikan $\rho \colon G \to S_4$. Tunjukkan bahwa
        $\ker \rho$, yang termuat di setiap $N_G(P_3)$ sehingga
        ordenya membagi $3$, bersifat trivial \emph{(mengapa ia tidak boleh
        berorde $3$?)}; bahwa petanya, yaitu subgrup berorde $12$ dari
        $S_4$, pastilah $A_4$ \emph{(subgrup berindeks 2 bersifat
        \omterm{def:b3:groups:normal}{normal} dan memuat semua kuadrat ---
        \cref{exo:b3:groups:1}; hitunglah kuadrat di $S_4$)}; lalu
        simpulkan $G \cong A_4$.
  \item Andaikan $n_3 = 1$, sehingga $G \cong \Z/3\Z \rtimes_\varphi P_2$
        dengan $\varphi \colon P_2 \to \operatorname{Aut}(\Z/3\Z)
        \cong \Z/2\Z$. Cacahlah semua kasusnya: $\varphi$ trivial
        melahirkan $\Z/12\Z$ dan $\Z/6\Z \times \Z/2\Z$; $P_2 =
        \Z/4\Z$ dengan $\varphi$ surjektif melahirkan
        \emph{grup disiklik} $\mathrm{Dic}_3 = \Z/3\Z \rtimes
        \Z/4\Z$; $P_2 = (\Z/2\Z)^2$ dengan $\varphi$ surjektif
        melahirkan, sampai kesetaraan pada pertanyaan 3, satu
        grup saja --- tunjukkan bahwa ia $D_6$, misalnya dengan menampilkan sebuah unsur
        berorde $6$ dan sebuah involusi bertipe pencerminan.
  \item Tunjukkan bahwa $\Z/12\Z$, $\Z/6\Z \times \Z/2\Z$, $D_6$, $A_4$,
        $\mathrm{Dic}_3$ tidak saling isomorfik \emph{(hitunglah
        unsur berorde $2$, atau pakailah $n_3$)}. Ini menuntaskan orde
        $12$.
\end{enumerate}

\medskip
\noindent\textbf{Bagian V --- Sintesis.}
\begin{enumerate}[resume]
  \item Susunlah tabel klasifikasinya: untuk setiap orde $n \leq
        15$, daftar lengkap grupnya sampai isomorfisma, dengan
        cacahan $1, 1, 1, 2, 1, 2, 1, 5, 2, 2, 1, 5, 1, 2, 1$.
\end{enumerate}

\medskip
\noindent\textbf{Bagian VI --- Lebih jauh: grup berorde $p^3$ dengan
$p$ ganjil.} Analisis orde $8$ pada Bagian III mempunyai padanan
yang indah untuk bilangan prima ganjil, dengan satu gejala yang benar-benar baru. Misalkan $p$
bilangan prima ganjil dan $G$ tak abelian berorde $p^3$.
\begin{enumerate}[resume]
  \item Tunjukkan bahwa $\abs{Z(G)} = p$, bahwa $G/Z(G) \cong
        (\Z/p\Z)^2$ \emph{(kuosien siklik terhadap pusat memaksa
        keabelian: \cref{exo:b3:groups:2})}, dan bahwa $D(G) =
        Z(G)$ \emph{(untuk $D(G) \subseteq Z(G)$, pakailah bahwa
        $G/Z(G)$ abelian; untuk kesamaannya, $G$ tak abelian dan
        $D(G) \neq \{e\}$)}. Simpulkan bahwa setiap \omterm{def:b3:groups:derived}{komutator}
        $[x, y] = xyx^{-1}y^{-1}$ bersifat pusat dan ordenya
        membagi $p$.
  \item (Identitas kuncinya) Misalkan $x, y \in G$ dan $z = [y, x]$,
        yang bersifat pusat. Buktikan dengan induksi pada $k$:
        \[
        (xy)^k = x^k y^k z^{k(k-1)/2} .
        \]
        \emph{(Geserlah setiap $y$ melewati setiap $x$; setiap penyeberangan berharga
        satu faktor pusat $z$.)}
  \item Simpulkan bahwa untuk $p$ ganjil pemetaan $\theta\colon x \mapsto
        x^p$ adalah morfisma grup dari $G$ ke $Z(G)$ \emph{(mengapa
        $x^p$ bersifat pusat? mengapa $z^{p(p-1)/2} = e$ memerlukan $p$
        ganjil?)}, lalu simpulkan bahwa eksponen $G$ adalah $p$ atau $p^2$,
        dan kedua kasus itu dibedakan oleh trivial tidaknya $\theta$.
  \item (Eksponen $p$) Andaikan setiap unsur memenuhi $x^p = e$.
        Pilih $x, y$ yang kelasnya membangun $G/Z(G)$ lalu tetapkan $z =
        [y, x]$. Tunjukkan bahwa $z \neq e$, bahwa setiap unsur $G$
        secara tunggal berbentuk $x^ay^bz^c$ ($0 \leq a, b, c < p$), dan bahwa
        perkaliannya sepenuhnya ditentukan oleh
        relasi $x^p = y^p = z^p = e$, $z$ bersifat pusat, $[y, x] =
        z$. Pastikan bahwa \emph{grup Heisenberg}
        \[
        H_p = \left\{
        \begin{pmatrix} 1 & a & c\\ 0 & 1 & b\\ 0 & 0 & 1
        \end{pmatrix} : a, b, c \in \mathbb F_p \right\}
        \subseteq GL_3(\mathbb F_p)
        \]
        mewujudkan relasi tersebut dan bereksponen $p$ (hitunglah
        $(I + N)^p$ dengan $N$ segitiga atas sejati, dengan memakai
        $N^3 = 0$ dan $p \geq 3$): setiap grup tak abelian bereksponen-$p$
        berorde $p^3$ isomorfik dengan $H_p$.
  \item (Eksponen $p^2$) Andaikan ada $r \in G$ yang berorde $p^2$,
        lalu tetapkan $N = \langle r\rangle$, yang \omterm{def:b3:groups:normal}{normal} (berindeks $p$:
        \cref{exo:b3:groups:10}). Tunjukkan bahwa ada $s \notin N$ dengan
        $s^p = e$ \emph{(ambil sembarang $t \notin N$; dengan memakai pertanyaan
        20, koreksilah ia: $\theta(t) = t^p \in Z(G) \subseteq N$
        --- berilah alasan bahwa $Z(G) = \langle r^p\rangle$ --- lalu pilih $a$
        sehingga $s = tr^{a}$ memenuhi $s^p = e$; di manakah keganjilan $p$
        dipakai?)}. Tunjukkan $srs^{-1} = r^{1+p}$ setelah $s$ diganti
        oleh suatu pangkatnya, lalu simpulkan: ada tepat \emph{satu}
        grup tak abelian berorde $p^3$ dan bereksponen $p^2$,
        yakni $\Z/p^2\Z \rtimes_\varphi \Z/p\Z$ dengan
        $\varphi(1)\colon r \mapsto r^{1+p}$ \emph{(pakailah pertanyaan
        3; $\operatorname{Aut}(\Z/p^2\Z)$ siklik berorde
        $p(p-1)$, diterima tanpa bukti di sini, sehingga ia mempunyai satu subgrup
        berorde $p$ saja)}.
  \item Simpulkan cacahannya: untuk $p$ ganjil ada tepat $5$
        grup berorde $p^3$ (tiga abelian, dua tak abelian),
        sama seperti untuk $p = 2$ --- tetapi kedua yang tak abelian
        bukan lagi $D_4$ dan $Q_8$. Tunjukkan dengan tepat di mana
        argumen untuk $p$ ganjil runtuh ketika $p = 2$: pada identitas di
        pertanyaan 19, $z^{k(k-1)/2}$ untuk $k = p = 2$ adalah $z^1 \neq
        e$, sehingga pengkuadratan bukan morfisma --- dan memang $Q_8$
        mempunyai satu unsur berorde $2$ saja, sedangkan yang
        bereksponen $4$, yaitu $D_4$, mempunyai lima.
\end{enumerate}

\medskip
\noindent\textbf{Bagian VII --- Pelengkap.}
\begin{enumerate}[resume]
  \item Untuk $p$ ganjil, hitunglah unsur berorde $p$ pada masing-masing
        dari kedua grup tak abelian berorde $p^3$: tunjukkan bahwa
        $H_p$ mempunyai tepat $p^3 - 1$ buah, sedangkan $M_p =
        \Z/p^2\Z \rtimes \Z/p\Z$ mempunyai tepat $p^2 - 1$
        \emph{(pakailah morfisma $\theta$ pada pertanyaan 20:
        kenali petanya, lalu orde kernelnya)}.
        Periksalah secara numerik untuk $p = 3$: $26$ lawan $8$.
        Jelaskan mengapa tidak ada argumen morfisma-pengkuadratan semacam ini
        yang dapat memisahkan $D_4$ dari $Q_8$, dan cacahan mana yang
        memisahkan keduanya.
  \item Sebuah bilangan bulat $n \geq 1$ disebut \emph{siklik} jika setiap grup
        berorde $n$ siklik. Tunjukkan bahwa jika $p^2 \mid n$ untuk
        suatu bilangan prima $p$, atau jika $n$ mempunyai pembagi prima $p < q$
        dengan $p \mid q - 1$, maka $n$ tidak siklik \emph{(pada
        masing-masing kasus tampilkan grup tak siklik berorde $n$, dengan memakai
        Bagian II untuk kasus kedua)}. Simpulkan bahwa $n$ siklik memaksa
        $\gcd(n, \varphi(n)) = 1$, dengan $\varphi$ adalah fungsi
        totien Euler, lalu cocokkan dengan tabel pada pertanyaan 17:
        di antara $n \leq 15$, orde yang hanya membawa satu grup
        tepat berupa $n \in \{1, 2, 3, 5, 7, 11, 13, 15\}$,
        persis yang memenuhi $\gcd(n, \varphi(n)) = 1$.

\end{enumerate}
\end{problem}
